% !TeX program = xelatex
% Model specification, revised 2026-10-27. Compile three times with XeLaTeX.
\documentclass[11pt,a4paper]{article}
\usepackage[margin=2.35cm,headheight=15pt]{geometry}
\usepackage{amsmath,amssymb,bm,mathtools}
\usepackage{booktabs,tabularx,longtable,array}
\usepackage[no-config]{fontspec}
\usepackage{needspace,enumitem,fancyhdr}
\usepackage[unicode,hidelinks]{hyperref}
\defaultfontfeatures{Ligatures=Common}
\defaultfontfeatures[\rmfamily,\sffamily]{Ligatures=Common,Mapping=}
\setmainfont[AutoFakeSlant=0.2]{Noto Serif CJK SC}
\setsansfont{Noto Sans CJK SC}
\setmonofont{Noto Sans Mono CJK SC}
\XeTeXlinebreaklocale "zh"
\XeTeXlinebreakskip=0pt plus 1pt minus 0.1pt
\setlength{\parindent}{2em}
\setlength{\parskip}{0.35em}
\setlength{\emergencystretch}{3em}
\setlength{\LTpre}{0.6em}
\setlength{\LTpost}{0.6em}
\renewcommand{\arraystretch}{1.22}
\setlist{itemsep=0.2em,topsep=0.4em,leftmargin=2em}
\allowdisplaybreaks[2]
\setcounter{tocdepth}{2}
\setcounter{secnumdepth}{3}
\numberwithin{equation}{section}
\renewcommand{\contentsname}{目录}
\renewcommand{\tablename}{表}
\pagestyle{fancy}
\fancyhf{}
\fancyhead[L]{\small Model 1024：规则信念与持续执行}
\fancyhead[R]{\small 2026-10-27}
\fancyfoot[C]{\thepage}
\newcolumntype{L}[1]{>{\raggedright\arraybackslash}p{#1}}
\newcolumntype{Y}{>{\raggedright\arraybackslash}X}
\newcommand{\ind}{\mathbb{I}}
\newcommand{\norm}[1]{\left\lVert#1\right\rVert}
\newcommand{\clip}{\operatorname{clip}}
\newcommand{\logit}{\operatorname{logit}}
\newcommand{\sigmoid}{\operatorname{sigmoid}}
\newcommand{\A}{\mathcal{A}}
\newcommand{\Hh}{\mathcal{H}}
\newcommand{\Rr}{\mathcal{R}}
\newcommand{\Dd}{\mathsf{D}}
\newcommand{\Ll}{\mathsf{L}}
\newcommand{\Ee}{\mathsf{E}}
\newcommand{\Ii}{\mathcal{I}}
\newcommand{\Tt}{\mathcal{T}}
\DeclareMathOperator*{\argmin}{arg\,min}
\DeclareMathOperator*{\argmax}{arg\,max}
\hypersetup{pdftitle={Model 1024：规则信念与持续执行的学习模型},pdfauthor={},bookmarksdepth=3}
\title{Model 1024\\[0.5em]
\Large 学习与决策闭环下的认知隐状态生成模型\\[0.5em]
\large 规则信念、持续执行与反馈学习}
\author{}
\date{2026年10月27日}
\begin{document}
\maketitle
\thispagestyle{empty}


Model 1024描述学习者如何根据反馈调整规则信念、探索新规则，并维持或改变当前执行的规则。
模型将当前保留并参与学习的规则称为活跃规则，以规则信念表示它们在当前集合中的相对可信度。
这些信念经过归一化，即各项均非负且总和为1。每题由一条执行规则产生选择；下一题的执行身份
由更新后的规则信念和延续当前执行方式的倾向共同决定。

规则信念与执行身份具有不同的作用。信念规定各规则的相对可信度，也用于选择局部搜索的起点
和确定容量已满时优先替换的规则。执行身份记录本题实际采用哪条规则。
一个被试在整段任务中具有固定的执行维持倾向，实际执行的规则则随试次变化。
这使模型可以表达正确规则已经在场、甚至获得较高信念，但作答仍由另一条规则产生的过程。

模型由决策与学习两个模块构成。决策模块形成内部刺激，并按当前执行规则的类别预测生成选择。
学习模块用实际反馈评价全部活跃规则，更新规则信念、判断精度和控制搜索的失败累积量，
随后完成搜索、容量替换、知识继承及下一题的执行规则选择。条件3同时学习类别配对关系。
条件1、2各有9项自由参数，条件3有10项。

本文依次说明模型的状态、逐试次机制、生成与推断、参数估计，以及三类错误来源和口述检验。
当前模型保留为研究候选：它能够表达持续执行及执行切换，执行错误的逐题恢复仍较弱。
模型对潜在过程的表达能力与行为资料对该过程的识别能力，分别通过生成和恢复检验评价。

\paragraph{阅读提示}
第\ref{sec:overview}节概括机制分工，第\ref{sec:mechanisms}节说明逐试次机制。
执行规则的维持与切换见第\ref{sec:execution}节；规则信念的保持与更新见第\ref{sec:retention}、\ref{sec:belief-update}节，
新规则的信念分配见第\ref{sec:admission}节，精度与配对继承见第\ref{sec:transfer}节。
参数定义与选参方法见第\ref{sec:inference}节，过程预测与检验范围见第\ref{sec:validation}节。

\paragraph{数据用途}
当前个体参数研究使用每名被试从首题到末题的全部有效选择，寻找同时解释实际选择和自主轨迹形态的参数组合。
自主生成通过模拟选择获得相应环境反馈，并形成完整学习路径。每个组合保留全部路径及各类形态的发生比例。
候选范围、数值精度和形态相容性共同限定可交付的参数集合；具体方案见第\ref{sec:selection}节。
本批真人记录均为开发资料，选择似然和生成检查参与个体选参，独立模拟种子用于确认计算稳定性。
此前前70\%选参的报告保留其原协议。口述资料仍待用户最终人工检查，新增口述验证与编码修改暂缓。

\clearpage
{\small\setlength{\parskip}{0pt}\tableofcontents}
\clearpage

\section{框架总览}
\label{sec:overview}

\subsection{两个心理模块与两个数学算子}
模型将一次学习试次分成决策与学习两个阶段。决策模块接收当前状态和刺激，产生选择；
学习模块结合该选择及反馈，更新下一题使用的状态：
\begin{equation}
(S_t,\bm x_t)\xrightarrow{\Dd_\theta}(m_t,y_t),\qquad
(S_t,m_t,y_t,r_t)\xrightarrow{\Ll_\theta}S_{t+1}.
\label{eq:loop}
\end{equation}
$S_t$是第$t$题作答前的认知状态，$\bm x_t$是物理刺激，$y_t$是选择，$r_t$是反馈。
$m_t$保存本题的内部知觉和规则预测，使学习使用与作答一致的信息。
$\theta$为同一被试在整段任务中保持固定的参数。

研究者在这个闭环上定义两个数学算子。\textbf{轨迹推断算子}$\Ii$结合实际行为，
更新各条可能内部路径的概率；\textbf{状态演化算子}$\Tt$对可能的选择和反馈求和，
给出模型自主学习的状态分布。粒子滤波是计算这些概率的数值方法，粒子代表研究者考虑的不同潜在路径。
认知容量$K$限制每条路径内的活跃规则数。

\subsection{参数、状态与试次记录}
参数规定学习过程的稳定属性，如搜索倾向、规则信念保持率、精度学习率和执行维持倾向。
状态随经验变化，记录当前考虑哪些规则、各规则的相对可信度、判断精度，以及本题的执行身份。
本题记录$m_t$连接决策和学习，在反馈更新后完成其作用。

心理含义有两个层次：参数的\emph{操作定义}由公式给出；参数能否可靠地描述某个被试，
取决于数据量、估计精度、参数补偿关系和恢复检验。
例如，失败敏感度为零时，搜索概率固定为基础搜索概率，失败累积量的保留参数在此边界上失去行为作用。
第\ref{sec:parameters}节将这些边界与参数定义放在一起。

\subsection{各认知机制的分工}
决策模块依次完成知觉与选择，学习模块完成反馈评价、精度学习、规则管理和下一题的执行准备。
条件3的反馈评价还包含对类别配对关系的学习。各机制通过当前状态和本题内部记录衔接。

{\small\noindent
\begin{tabularx}{\linewidth}{L{2.5cm}L{3.3cm}Y}
\toprule
所属阶段 & 机制 & 输入、输出和作用\\
\midrule
决策 & 知觉 & 根据物理刺激和个体知觉误差形成内部刺激。\\
决策 & 选择 & 按当前执行规则的类别预测生成选择概率。\\
学习 & 反馈评价 & 计算反馈似然，调节并更新规则信念；条件3同时更新配对信念。\\
学习 & 精度学习 & 根据对数反馈似然的梯度调整各活跃规则的判断精度。\\
学习 & 规则管理 & 更新失败累积量，搜索新规则，完成容量替换、入场信念分配与知识继承。\\
学习 & 执行准备 & 根据新集合的规则信念及上一执行身份，抽取下一题的执行规则。\\
\bottomrule
\end{tabularx}}


\subsection{规则在场、规则信念与执行身份}
规则在场表示这条规则仍保留在活跃集合中，能够继续接受反馈检验。
规则信念表示它在当前集合中的相对可信度。执行身份表示本题由哪条规则产生回答。
这三个量分别描述可用的判断、对判断的评价和实际采用的判断。

每次反馈先调节已有信念的保持程度，再通过反馈似然改变各规则的相对比例。
反馈似然是规则预测当前实际反馈的概率，衡量它对本次经验的解释程度。
信念较高的规则更可能成为局部搜索的起点；新规则进入已满集合时，信念最低的已有规则被替换。
新集合内的信念经过入场分配后仍然归一化。

下一题的执行规则根据信念和当前执行身份抽取。执行维持倾向提高继续使用同一条规则的概率，
因此信念的变化与执行的变化可以发生在不同试次。
判断精度决定执行规则如何把类别距离转化为选择概率；全部活跃规则继续进行反馈学习。
失败累积量控制探索时机，知识继承使已形成的判断精度和配对知识延续到新规则。

\section{实验设计、规则空间与认知状态}
\label{sec:state}

\subsection{实验条件与观测}
第$t$个有效试次包含四维归一化长度、实际选择和实际反馈：
\begin{equation}
\bm x_t\in[0,1]^4,\qquad y_t\in\{1,\ldots,C\},\qquad r_t\in\Rr.
\end{equation}
四维对应实验中的四个部位，在整段任务中保持相同顺序。条件1为二分类二值反馈，
$C=2$、$\Rr=\{0,1\}$；条件2为四分类二值反馈，$C=4$、$\Rr=\{0,1\}$；
条件3为四分类三值反馈，$C=4$、$\Rr=\{0,0.5,1\}$。
条件3选对种类得1分，同科但种类错误得0.5分，另一科得0分。

任务输入$u_t$包含刺激、任务条件及环境反馈分布$\Ee_{u_t}(r\mid\bm x_t,y)$。
真实类别与按键映射用于环境评分及数据坐标转换；学习者根据自己的选择和获得的反馈学习。
条件3的真实类别、统一反应编号和实际按键分别记录，未知配对从均匀信念开始。
漏答试次跳过，模型按有效试次的原始顺序递推，并保留其原始试次编号。

\subsection{规则目录与活跃集合}
规则目录为$\Hh=\{1,\ldots,J\}$，二分类共有29条，四分类共有152条，分组见附录\ref{app:catalogue}。
每条规则规定一组判断、类别区域及固定类别标签。当前活跃集合$\A_t$由容量参数$K$约束：
\begin{equation}
\mu(h)=1/J,\qquad \A_t\subset\Hh,\qquad 1\leq|\A_t|\leq K<J.
\label{eq:capacity}
\end{equation}
$\mu$为基础规则权重，当前估参允许$K\in\{2,3,4,6,8\}$。
模型从目录中发现规则，目录的几何形式与标签共同限定其可表达策略。


\subsection{决策就绪状态与初始化}
把前题反馈引起的学习、规则管理和执行准备完成后，本题呈现刺激前的状态写成
\begin{equation}
S_t=\left(\A_t,e_t,\{\pi_t(h),\beta_t(h)\}_{h\in\A_t},F_t,
\{q_t(k\mid h)\}_{h\in\A_t,\ k=1,2,3}\right).
\label{eq:state}
\end{equation}
$e_t\in\A_t$为本题执行规则的身份，$\pi_t(h)$为规则信念，
$\beta_t(h)$为判断精度，$F_t$为失败累积量。
规则信念满足$\pi_t(h)\geq0$、$\sum_{h\in\A_t}\pi_t(h)=1$。
$q_t(k\mid h)$为给定规则$h$的条件配对信念，仅条件3使用，且每条规则的三种配对概率之和为1。
条件3也可用联合分布$\Pi_t(h,k)=\pi_t(h)q_t(k\mid h)$表示这两层信念。

首题抽取$h_1\sim\mu$，设置$\A_1=\{h_1\}$、$\pi_1(h_1)=1$、
$\beta_1(h_1)=\beta_0$、$F_1=0$、$e_1=h_1$；条件3令$q_1(k\mid h_1)=1/3$。
首题完成后才根据其反馈启动搜索控制。

本题内部记录为
\begin{equation}
m_t=\left(\widetilde{\bm x}_t,e_t,\{d_{t,h,c}\},\{p_{t,h}(c)\},\{\pi_t(h)\}\right),
\end{equation}
分别保存内部刺激、执行身份、类别距离、规则类别预测和作答前信念。
学习模块使用这份记录与本题反馈完成更新，形成下一有效试次的$S_{t+1}$。
$S_t$与$S_{t+1}$对应相同的时点：前一试次学习已完成，当前试次尚未作答。

\section{认知机制设计}
\label{sec:mechanisms}

\subsection{决策模块：从当前状态生成选择}
\label{sec:decision}

\subsubsection{知觉：从物理刺激形成内部刺激}
决策模块的输入是$S_t$与$\bm x_t$，输出本题内部记录和选择，其概率规律记为
$\Dd_\theta(dm,y\mid S,\bm x)$。它先形成内部刺激，再计算类别预测，并按当前执行规则生成选择。
各规则共享同一次知觉结果：
\begin{equation}
\widetilde x_{t,j}=\clip_{[0,1]}(x_{t,j}+\epsilon_{t,j}),\qquad j=1,\ldots,4.
\label{eq:perception}
\end{equation}
$\clip$将越界值限制在指定区间，$\epsilon_{t,j}$表示第$j$个特征的知觉误差。
误差采用正态或均匀分布，其分布形式与参数由独立知觉测量确定，并与个体和特征对应。
行为参数估计期间，这些测量参数保持固定。

\subsubsection{类别预测：从区域距离计算类别概率}
规则$h$将刺激空间划分为类别区域$\mathcal R_{h,c}$。
模型用内部刺激到各区域的最短欧氏距离计算类别概率：
\begin{align}
d_{t,h,c}&=\inf_{\bm z\in\mathcal R_{h,c}}\norm{\widetilde{\bm x}_t-\bm z}_2,\\
p_{t,h}(c)&=\frac{\exp[-\beta_t(h)d_{t,h,c}]}
{\sum_{a=1}^C\exp[-\beta_t(h)d_{t,h,a}]}.
\label{eq:category}
\end{align}
刺激位于某区域内时距离为0；类别有多个分离分量时，取最近分量的距离。
$\beta$越大，概率越集中在距离较小的类别。知觉误差改变内部刺激，
精度改变类别距离对判断的影响，两者作用于不同环节。


\subsubsection{单规则选择：按执行规则的类别预测作答}
当前执行身份$e_t$已在本题刺激呈现前确定。给定状态和内部刺激，本题选择概率为
\begin{equation}
p_t(c\mid S_t,\widetilde{\bm x}_t)=p_{t,e_t}(c).
\label{eq:choice}
\end{equation}
右侧是执行规则对类别$c$的预测，全部类别的概率之和为1。
自主生成从这一分布抽取$y_t$；行为计分读取实际选择的概率。
执行规则的预测本身具有概率性，因此执行正确规则时仍可能产生错误选择。

例如，两条规则的信念分别为0.8和0.2，当前执行第二条规则。
若第二条规则给类别1的概率为0.3，本题选择类别1的概率就是0.3。
两条规则的信念将在学习和下一次执行选择中发挥作用。
下一执行身份的抽取规则见第\ref{sec:execution}节。

一条潜在路径内，每题具有一个明确的执行身份。研究者无法直接观察该身份，
因此推断时对多条可能路径的选择概率加权平均，得到行为预测。
这种跨路径平均表达研究者对内部状态的不确定性；每条路径仍按式\eqref{eq:choice}作答。
额外随机作答率固定为0，行为变异来自知觉、规则搜索、执行切换和类别概率抽样。

\subsection{学习模块：从反馈形成下一状态}
\label{sec:learning}

\subsubsection{学习顺序：用本题预测评价反馈}
学习模块接收$S_t,m_t,y_t,r_t,u_t$，其输出分布记为
$\Ll_\theta(dS'\mid S,m,y,r,u)$。
它用作答时的类别预测计算反馈似然，再更新规则信念、精度与失败累积量，
随后提出候选、处理容量、初始化新规则，最后选定下一题的执行规则。条件3同时更新配对知识。

\textbf{反馈评价使用作答时的预测。}信念更新与精度梯度共享$m_t$中的知觉、距离和类别概率。
精度学习后的值用于后续判断和新规则继承。
因此，反馈先评价本题作答时的预测，再改变后续试次的判断。

\subsubsection{二值任务的反馈似然}
反馈似然$\ell_t(h)$衡量规则$h$对当前实际反馈的预测程度。
正确反馈对应所选类别，错误反馈对应其余类别的总概率：
\begin{equation}
\ell_t(h)=
\begin{cases}
p_{t,h}(y_t),&r_t=1,\\
1-p_{t,h}(y_t),&r_t=0.
\end{cases}
\label{eq:binary-feedback}
\end{equation}
因此，四分类中的错误反馈评价三个其他类别的总概率。
所有活跃规则都接收这次反馈评价，各规则的反馈似然由其自己的预测决定。

\subsubsection{规则信念保持：根据反馈调整已有判断}
\label{sec:retention}
模型假设，成功的判断使学习者更充分地保持已有信念，失败则降低对已有判断的依赖。
设$\delta\in[0,1]$为个体的基础规则信念保持率，本题有效保持率为
\begin{equation}
\delta_t=\delta+(1-\delta)r_t.
\label{eq:effective-retention}
\end{equation}
正确反馈$r_t=1$时，$\delta_t=1$；错误反馈$r_t=0$时，$\delta_t=\delta$；
半分反馈$r_t=0.5$时，$\delta_t=(1+\delta)/2$。
例如$\delta=0.8$时，三种反馈对应的有效保持率依次为1、0.8和0.9。

在结合本题反馈似然之前，先对已有信念作幂变换并归一化：
\begin{equation}
\bar\pi_t(h)=\frac{\pi_t(h)^{\delta_t}}
{\sum_{v\in\A_t}\pi_t(v)^{\delta_t}}.
\label{eq:belief-retention}
\end{equation}
当$\delta_t=1$时保留原有比例；当$0<\delta_t<1$时，正信念之间的比值更接近1。
当$\delta_t=0$时，约定全部活跃规则的$\bar\pi_t$均为$1/|\A_t|$，包括此前信念为零的规则。
这一步完成后，再由反馈似然决定哪些规则应获得更高的更新后信念。

$\bar\pi_t$是收到反馈后用于学习的中间分布。本题选择已由执行规则$e_t$的类别预测生成，
当前反馈通过$\delta_t$和后续似然更新共同改变下一题的状态。
这是一项反馈依赖的信念学习假设，其心理解释需要结合行为与内部状态检验。

\subsubsection{条件3的配对信念与反馈似然}
\label{sec:pairing}
条件3的四个反应有三种可能的两两配对：
\begin{equation}
G_1=\{\{1,2\},\{3,4\}\},\quad G_2=\{\{1,3\},\{2,4\}\},\quad
G_3=\{\{1,4\},\{2,3\}\}.
\end{equation}
令$\operatorname{mate}_k(y)$为配对$G_k$中与$y$同科的另一个反应。
对于规则$h$和配对$k$，本题反馈似然为
\begin{equation}
\ell_t(h,k)=\begin{cases}
p_{t,h}(y_t),&r_t=1,\\
p_{t,h}(\operatorname{mate}_k(y_t)),&r_t=0.5,\\
1-p_{t,h}(y_t)-p_{t,h}(\operatorname{mate}_k(y_t)),&r_t=0.
\end{cases}
\label{eq:pair-likelihood}
\end{equation}
满分、半分和零分分别对应所选种类、同科另一种类和另一科的两个种类。


条件配对信念由保持率$\gamma$调整：
\begin{equation}
\bar q_t(k\mid h)=\frac{q_t(k\mid h)^\gamma}
{\sum_{j=1}^3q_t(j\mid h)^\gamma}.
\label{eq:pair-forgetting}
\end{equation}
$\gamma=0$时，每条规则的三种配对设为均匀分布；$\gamma=1$时完整保留原比例。
规则层使用式\eqref{eq:belief-retention}的$\bar\pi_t$，条件配对层使用$\bar q_t$，
两层分别归一化。规则对实际反馈的预测概率为
\begin{equation}
\ell_t(h)=\sum_k\bar q_t(k\mid h)\ell_t(h,k).
\label{eq:marginal-likelihood}
\end{equation}
模型据此形成规则与配对的联合更新：
\begin{align}
\Pi_t^+(h,k)&=
\frac{\bar\pi_t(h)\bar q_t(k\mid h)\ell_t(h,k)}
{\sum_{v\in\A_t}\sum_j\bar\pi_t(v)\bar q_t(j\mid v)\ell_t(v,j)},\\
\pi_t^+(h)&=\sum_k\Pi_t^+(h,k),\qquad
q_t^+(k\mid h)=\frac{\Pi_t^+(h,k)}{\pi_t^+(h)},\qquad
Q_t^+(k)=\sum_h\Pi_t^+(h,k).
\label{eq:pair-update}
\end{align}
这里$\Pi_t^+$是反馈后的联合信念，$\pi_t^+$是规则边际，$q_t^+$是各规则的条件配对信念。
规则边际为零时，约定其条件配对分布为均匀分布。
$Q_t^+$汇总全部已有规则的配对知识，供新规则继承。
精度学习使用的配对权重为$\bar q_t$，与本题反馈似然的计算时点一致。

两项保持率的分工是：$\delta_t$调节规则之间的信念比例，$\gamma$调节各规则内部的配对比例。
两层再通过同一反馈似然形成联合信念。条件1、2使用式\eqref{eq:binary-feedback}，
配对保持率$\gamma$仅在条件3估计。

\subsubsection{规则信念更新：结合已有判断与反馈似然}
\label{sec:belief-update}
三种条件均用本题反馈似然重新分配规则信念：
\begin{equation}
\pi_t^+(h)=\frac{\bar\pi_t(h)\ell_t(h)}
{\sum_{v\in\A_t}\bar\pi_t(v)\ell_t(v)}
=\frac{\pi_t(h)^{\delta_t}\ell_t(h)}
{\sum_{v\in\A_t}\pi_t(v)^{\delta_t}\ell_t(v)}.
\label{eq:belief-update}
\end{equation}
指数为零时按第\ref{sec:retention}节的均匀约定处理。
条件3的式\eqref{eq:belief-update}正是式\eqref{eq:pair-update}的规则边际，
因此每次反馈只形成一次规则信念更新。

在保持后信念相同时，反馈似然越大的规则越能解释实际反馈，获得的更新后信念也越大。
例如，保持后的两条规则信念为0.6和0.4，反馈似然分别为0.8和0.2，
则未归一化质量为0.48和0.08，更新后信念约为0.857和0.143。
两条规则均使用本次反馈；更新结果同时取决于已有比例和预测反馈的能力。

对于信念与似然均为正的两条规则$h,v$，更新还可写成对数比值的形式：
\begin{equation}
\log\frac{\pi_t^+(h)}{\pi_t^+(v)}
=\delta_t\log\frac{\pi_t(h)}{\pi_t(v)}
+\log\frac{\ell_t(h)}{\ell_t(v)}.
\label{eq:belief-odds}
\end{equation}
第一项保留此前的相对判断，第二项加入当前反馈对两条规则的区分。
此前经验通过当前信念持续影响学习，其影响随各题的有效保持率及规则集合变化而调整。
更新后的$\pi_t^+$用于搜索、容量替换和继承，随后经新规则入场分配形成$\pi_{t+1}$。

\subsubsection{精度学习：按反馈调整类别区分程度}
\label{sec:precision}
精度学习使用对数反馈似然对$\beta$的梯度，表示小幅改变精度时，对数反馈似然如何变化。
令$a_{t,h,c}$为类别$c$与实际反馈相容的程度。二值任务中，
$r=1$时$a_{t,h,c}=\ind[c=y_t]$，$r=0$时$a_{t,h,c}=\ind[c\ne y_t]$。
条件3对三种配对加权：
\begin{equation}
a_{t,h,c}=\sum_k\bar q_t(k\mid h)\ind[c\in Y_k(y_t,r_t)],
\end{equation}
其中，$\ind$为指示函数，条件成立时取1，否则取0；$Y_k$是式\eqref{eq:pair-likelihood}中与实际反馈对应的类别集合。
将本题类别预测按反馈相容程度加权并归一化，得到反馈条件下的类别分布：
\begin{equation}
p_{t,h}^{\mathrm{fb}}(c)=\frac{a_{t,h,c}p_{t,h}(c)}{\ell_t(h)}.
\end{equation}
固定本题知觉、距离和反馈前配对分布，梯度可写成两个平均距离之差：
\begin{equation}
g_t^\beta(h)=\left.\frac{\partial\log\ell_t(h;\beta)}{\partial\beta}\right|_{\beta_t(h)}
=\sum_c p_{t,h}(c)d_{t,h,c}-\sum_c p_{t,h}^{\mathrm{fb}}(c)d_{t,h,c}.
\label{eq:gradient}
\end{equation}
反馈相容类别的平均距离较小时，提高精度；较大时，降低精度。
更新方向由反馈与规则的实际预测共同确定。

判断精度按梯度作加法更新，并限制在$[0.1,25]$内：
\begin{equation}
\beta_t^+(h)=\clip_{[0.1,25]}\left[\beta_t(h)+10\eta g_t^\beta(h)\right].
\label{eq:beta}
\end{equation}
固定系数10规定学习率的数值尺度，$\eta$为个体参数，有效步长为$10\eta$。
每条活跃规则均按自己的梯度更新。有限步长和截断共同决定实际变化，
单步似然的实际增减需由该更新结果计算。

\subsubsection{失败累积量与搜索触发}
失败累积量$F_t$控制下一题前是否探索新规则。完全成功使已有失败量衰减，
不完全成功则按失分程度增加：
\begin{equation}
F_{t+1}=\begin{cases}\lambda F_t,&r_t=1,\\F_t+(1-r_t),&r_t<1.\end{cases}
\label{eq:failure}
\end{equation}
因此条件3的半分增加0.5，零分增加1。$\lambda\in[0,1]$表示一次完全成功之后保留多少失败量。
搜索概率为
\begin{equation}
s_{t+1}=\sigmoid\!\left[\logit(p_0)+a_eF_{t+1}\right],\qquad
\sigmoid(z)=\frac1{1+e^{-z}},\quad\logit(p)=\log\frac{p}{1-p}.
\label{eq:search}
\end{equation}
$p_0\in(0,1)$为失败量为零时的基础搜索概率，$a_e\geq0$为失败敏感度。
每次抽取一个概率为$s_{t+1}$的搜索事件，发生时提出一条新规则。
完全成功后的基础搜索机会允许模型继续探索。

\subsubsection{局部结构搜索与全局搜索}
\label{sec:proposal}
候选范围为$\mathcal U_t=\Hh\setminus\A_t$。
搜索发生后，以固定个体比例$g_0$采用全局路线，从$\mu$在$\mathcal U_t$上的归一化分布提出候选。

其余情况采用局部路线，按反馈更新后的规则信念抽取父规则$h_{\mathrm{parent}}$：
\begin{equation}
\Pr(h_{\mathrm{parent}}=h\mid\text{局部搜索},S_t,m_t,y_t,r_t)=\pi_t^+(h).
\label{eq:parent}
\end{equation}
信念较高的规则更可能成为进一步修改的起点。

结构邻域$\mathcal N_t(h)\subset\mathcal U_t$包含对规则$h$作一次合法结构修改后得到的规则。
编辑距离$d_{\mathrm{edit}}(h,u)$记录到候选$u$的最小单步成本；更换特征、
改变比较关系及标准参照、修改参照值、改写条件树一条下级判断的成本依次为1、1.5、1、2。
编辑保留其余判断和相应的类别输出布局，且结果必须在固定目录内。
非空邻域中的提议概率为
\begin{equation}
Q_{\mathrm{edit},t}(u\mid h)=
\frac{\ind[u\in\mathcal N_t(h)]e^{-d_{\mathrm{edit}}(h,u)}}
{\sum_{v\in\mathcal N_t(h)}e^{-d_{\mathrm{edit}}(h,v)}}.
\end{equation}
成本越低的修改越容易提出。

当结构邻域没有非活跃规则时，局部搜索依据固定类别标签下的功能相似性提出候选。
功能距离$d_{\mathrm{func}}(h,u)$是两条规则在均匀参考刺激上给出不同类别标签的比例。
该距离以100000个共同参考刺激估计，随机种子为0。功能搜索的提议分布为
\begin{equation}
Q_{\mathrm{func},t}(u\mid h)=
\frac{\ind[u\in\mathcal U_t]\mu(u)e^{-d_{\mathrm{func}}(h,u)/0.10}}
{\sum_{v\in\mathcal U_t}\mu(v)e^{-d_{\mathrm{func}}(h,v)/0.10}}.
\end{equation}
记$Q_{\mathrm{local},t}$为有邻居时的结构分布、无邻居时的功能回退分布，
则搜索事件内的完整候选分布为
\begin{equation}
Q_t(u)=g_0\frac{\mu(u)}{\sum_{v\in\mathcal U_t}\mu(v)}
+(1-g_0)\sum_{h\in\A_t}\pi_t^+(h)Q_{\mathrm{local},t}(u\mid h),
\quad u\in\mathcal U_t.
\label{eq:proposal}
\end{equation}
新候选的概率由当前活跃状态和固定规则关系确定。
局部搜索在邻域非空时按结构编辑距离抽样，邻域为空时按功能距离抽样。

\subsubsection{容量替换与规则保留}
\label{sec:replacement}
没有搜索事件时，$\A_{t+1}=\A_t$。
发生搜索并提出$u$后，有空位就加入；容量已满就从已有规则中移除信念最低的一条：
\begin{equation}
\A_{t+1}=\begin{cases}
\A_t\cup\{u\},&|\A_t|<K,\\
(\A_t\setminus\{h^-\})\cup\{u\},&|\A_t|=K,
\end{cases}
\qquad h^-\in\argmin_{h\in\A_t}\pi_t^+(h).
\label{eq:replacement}
\end{equation}
多条规则的信念并列最低时，等概率抽取其中一条。
比较对象为已有规则，新候选在搜索发生时获得一次入集机会。

因此，规则保留取决于当前集合是否出现新候选，以及它在已有规则中的相对信念。
信念较低的规则可继续留在尚未搜索的集合中，并可能因执行维持倾向而继续产生回答。
当前执行规则适用同一替换规定；它被移除后，下一题必须采用其他规则。
规则集合从首题的一条逐渐增加，到达容量后通过替换改变组成。


\subsubsection{新规则入场：按反馈分配初始信念}
\label{sec:admission}
新规则需要获得初始信念，才能获得后续执行和局部搜索的选择概率。模型根据集合是否已满，以及刚完成试次的
反馈，为新规则分配概率质量。没有搜索事件时，直接令$\pi_{t+1}=\pi_t^+$。

发生搜索并提出新规则$u$时，记其入场份额为$\nu_{t+1}$。
当活跃集合尚未达到容量时，设$n_t=|\A_t|$，新规则获得新集合的均匀份额$1/(n_t+1)$。
当集合已满时，设$h^-$为本次被移除的规则，其份额按反馈调节：
\begin{equation}
\nu_{t+1}=\begin{cases}
\displaystyle\frac{1}{n_t+1},&n_t<K,\\[0.5em]
\displaystyle r_t\pi_t^+(h^-)+\frac{1-r_t}{K},&n_t=K.
\end{cases}
\label{eq:entry-share}
\end{equation}
满容量时，答对后的新规则接替被移除规则的份额；答错后获得$1/K$的试用份额；
半分反馈取二者的平均。例如$K=4$且被移除规则的信念为0.001时，正确、半分和错误
反馈后的新规则份额依次为0.001、0.1255和0.25。该分配表示判断失败时更充分地尝试新解释。

记幸存规则集合为$\A_t^{\mathrm{keep}}=\A_t\cap\A_{t+1}$。
新规则获得$\nu_{t+1}$，幸存规则按更新后的相对比例分配剩余质量：
\begin{equation}
\pi_{t+1}(u)=\nu_{t+1},\qquad
\pi_{t+1}(h)=(1-\nu_{t+1})
\frac{\pi_t^+(h)}{\sum_{v\in\A_t^{\mathrm{keep}}}\pi_t^+(v)},
\quad h\in\A_t^{\mathrm{keep}}.
\label{eq:entry-belief}
\end{equation}
这一步保持幸存规则之间的信念比值，并使新集合内的总信念为1。
满容量且答对时，幸存规则的绝对信念也保持原值。新规则从下一题开始接受自身的反馈检验。
式\eqref{eq:entry-share}中的$r_t$始终是新规则入场之前刚完成试次的反馈。

被移除规则的信念至多为$1/K$，因此错误后的均匀试用份额至少与其原份额相同。
若被移除规则的信念恰为零，答对后的新规则也得到零信念；有效保持率大于零时，
后续乘法更新保留该零值。有效保持率为零时，活跃规则重新获得均匀的保持后分布。
这些边界说明，规则在场、规则信念与执行身份需要分别检查。

\subsubsection{精度与配对继承：延续已有任务知识}
\label{sec:transfer}
保留规则延续本题学习后的判断精度，新规则继承全部已有活跃规则的信念加权平均精度：
\begin{equation}
\bar\beta_t=\sum_{h\in\A_t}\pi_t^+(h)\beta_t^+(h),\qquad
\beta_{t+1}(h)=\begin{cases}\beta_t^+(h),&h\in\A_t^{\mathrm{keep}},\\
\bar\beta_t,&h=u.\end{cases}
\label{eq:transfer}
\end{equation}
参与平均的是\textbf{容量替换之前的全部活跃规则}，包括本次随后被移除的规则；
信念与精度均已吸收本题反馈。信念非负且和为1，因此继承值位于参与平均的最小精度与最大精度之间。
首题规则使用$\beta_0$。

这一继承假设任务中形成的刺激区分能力能够在规则之间延续。
例如，已有规则的更新后精度为6和10、信念为0.75和0.25时，新规则精度为7。
其后每条规则仍按自己的几何预测和反馈调整精度。

条件3中，保留规则延续自己的条件配对信念，新规则继承式\eqref{eq:pair-update}的整体配对边际：
\begin{equation}
q_{t+1}(k\mid h)=\begin{cases}q_t^+(k\mid h),&h\in\A_t^{\mathrm{keep}},\\
Q_t^+(k),&h=u,\end{cases}\qquad
\Pi_{t+1}(h,k)=\pi_{t+1}(h)q_{t+1}(k\mid h).
\label{eq:pair-transfer}
\end{equation}
$Q_t^+$同样在容量替换前计算，以更新后的规则信念汇总各规则的配对知识。
局部与全局提出的新规则采用相同的信念分配及继承方式。

规则离开活跃集合时，其专属信念、精度与条件配对状态随之结束。
将来重新提出时，按届时的反馈与活跃状态分配入场信念，并继承精度及配对知识。

\subsubsection{执行准备：根据信念维持或切换规则}
\label{sec:execution}
完成反馈学习、规则管理和知识继承后，模型从新活跃集合中选择下一题的执行规则。
令$\kappa_s\geq0$表示被试$s$的执行维持倾向，在一次拟合的整段任务中固定。
以下省略被试下标，用$\omega_{t+1}(h)$表示选择规则$h$作为下一执行身份的概率：
\begin{equation}
\omega_{t+1}(h)=
\Pr(e_{t+1}=h\mid e_t,\pi_{t+1},\A_{t+1})=
\frac{\pi_{t+1}(h)\exp\!\left[\kappa\ind(h=e_t)\right]}
{\sum_{v\in\A_{t+1}}\pi_{t+1}(v)\exp\!\left[\kappa\ind(v=e_t)\right]},
\quad h\in\A_{t+1}.
\label{eq:execution}
\end{equation}
指示函数$\ind(h=e_t)$在候选与本题执行规则相同时取1，其余情况取0。
因此，当前执行规则的信念乘以$e^\kappa$后，与其他规则的信念一起归一化。
信念规定规则的相对可信度，执行维持倾向增加延续现有作答方式的概率。
抽取一次$e_{t+1}$后，它成为下一题状态的一部分。

当当前执行规则仍在集合中时，记$b=\pi_{t+1}(e_t)$，继续执行它的概率为
\begin{equation}
\Pr(e_{t+1}=e_t)=\frac{b e^\kappa}{b e^\kappa+1-b}.
\label{eq:execution-stay}
\end{equation}
例如，两条规则的信念为0.2和0.8，当前执行前者且$\kappa=2$，继续执行它的概率约为0.65。
如果切换到信念为0.8的规则，且两条规则的信念暂时保持不变，继续执行新规则的概率约为0.97。
这说明较高的规则信念可以先于稳定执行出现；执行切换也可以带来较突然的行为变化。
实际任务中，信念和活跃集合随反馈变化，持续执行的长度由整段状态演化共同决定。

执行规则被容量机制移除时，式\eqref{eq:execution}中的所有指示函数均为0，
下一执行身份直接按新集合的信念抽取。模型分别记录这种强制切换，
以及原规则仍在场时发生的自主切换。较大的$\kappa$提高维持概率，
其实际作用仍取决于当前规则的信念和它能否继续留在活跃集合中。

当$\kappa=0$时，$\omega_{t+1}=\pi_{t+1}$，每题按当前信念重新抽取执行身份。
在这一边界上，对执行身份求和后的类别概率为$\sum_h\pi_t(h)p_{t,h}(c)$。
全部活跃规则继续学习，使该边界的完整行为分布也由信念加权预测确定。
执行身份的逐题推断在此边界上需要结合行为等价关系解释。

执行准备发生在下一题刺激呈现之前，因此本题反馈只影响后续执行。
模型中被试层面的固定量是$\kappa_s$，试次层面变化的状态是$e_t$。
两者分别规定延续执行的倾向与每题实际采用的规则。

\subsubsection{一次完整有效试次}
\begin{enumerate}
\item 从$S_t$读取活跃规则、规则信念、精度、配对信念及执行身份$e_t$，抽取本题内部刺激。
\item 计算各活跃规则的类别预测，按执行规则的预测生成选择或计算实际选择的概率。
\item 接收实际反馈；自主生成时，由环境对模拟选择评分。
\item 根据$r_t$确定有效保持率$\delta_t$；用作答时的预测计算反馈似然，更新全部活跃规则的信念。
条件3同时调整条件配对信念并完成一次联合更新。
\item 用同一份作答时预测更新各规则的判断精度，更新失败累积量，形成下一题前的搜索概率。
\item 搜索发生时按局部／全局路线提出新规则；容量已满时替换信念最低的已有规则。
\item 根据本题反馈为新规则分配入场信念，按更新后的信念继承平均精度；条件3继承综合配对知识。
\item 根据新集合的信念、$e_t$和$\kappa$抽取$e_{t+1}$，记录执行维持或切换。
\item 将新集合及其信念、精度、配对信念、失败累积量和执行身份组成$S_{t+1}$。
\end{enumerate}
反馈信念更新与精度学习使用原有活跃集合。新规则的初始状态由已有知识与入场分配确定，
它可在下一题开始执行，其自身的首次反馈学习发生在下一有效试次。

\section{试次级隐状态生成}
\label{sec:kernel}

\subsection{联合生成核}
决策、环境反馈与学习共同规定从当前状态出发，下一状态及观测如何联合产生：
\begin{equation}
k_\theta(dS',y,r\mid S,u)=
\int \Dd_\theta(dm,y\mid S,\bm x)\,
\Ee_u(r\mid\bm x,y)\,
\Ll_\theta(dS'\mid S,m,y,r,u).
\label{eq:kernel}
\end{equation}
这里的积分消去内部试次记录 $m$，并保留该记录在决策与学习之间的依赖关系。
核使用概率测度记法，以同时涵盖连续的知觉和精度、离散的规则集合及确定性的状态更新。
固定 $S,u$ 后，对 $S'$ 积分并对全部 $y,r$ 求和的结果为1。

\subsection{生成、计分与内部推断}
自主生成从 $\Dd$ 抽取选择，由环境产生相应反馈，再经 $\Ll$ 更新状态。
观测计分固定实际选择与反馈，对可能的知觉、规则集合和其他状态进行积分。
内部推断则以这些观测重新分配潜在路径的概率。三者共享同一组参数与更新机制。

当实验评分确定时，对有效记录有 $\Ee_u(r_t\mid\bm x_t,y_t)=1$。
因此研究者的行为似然由选择概率构成，反馈作为已给定的学习输入改变下一状态。
规则的反馈似然 $\ell_t(h)$ 用于更新规则信念、判断精度和配对信念；粒子路径的观测权重按外显选择更新。
这一区分保证同一反馈的认知更新与行为序列计分各有明确位置。

\section{学习轨迹推断与状态分布演化}
\label{sec:operators}

\subsection{轨迹推断算子}
令 $b$ 为当前决策就绪状态的概率分布，$A'$ 为下一状态空间中的可测集合。
给定实际观测 $o=(y,r)$，定义未归一化算子
\begin{equation}
(\Ii_{\theta,o}^{u}b)(A')=\int k_\theta(A',o\mid S,u)\,b(dS).
\label{eq:inference}
\end{equation}
它按实际观测的生成概率重加权路径，并把路径推进到学习后的状态。
记 $\langle1,\zeta\rangle$ 为测度 $\zeta$ 的总质量，则
\begin{equation}
P_\theta(o\mid u,b)=\langle1,\Ii_{\theta,o}^{u}b\rangle,\qquad
b_{t+1}=\frac{\Ii_{\theta,o_t}^{u_t}b_t}
{\langle1,\Ii_{\theta,o_t}^{u_t}b_t\rangle}.
\label{eq:filter}
\end{equation}
第一式给出观测概率，第二式给出观察本题并完成学习后的状态后验。
未归一化的 $\Ii$ 对输入测度是正线性的；除以总质量的归一化过滤映射一般为非线性。

对给定输入序列，完整观测序列概率可写为
\begin{equation}
P_\theta(o_{1:T}\mid u_{1:T})=
\left\langle1,\Ii_{\theta,o_T}^{u_T}\cdots\Ii_{\theta,o_1}^{u_1}b_1\right\rangle.
\label{eq:sequence-operator}
\end{equation}
算子从右向左按时间作用。式\eqref{eq:filter}给出各时点的过滤边际；
完整路径的联合后验还需保留路径间的祖先关系。
利用未来观测更新过去状态属于平滑推断。逐试次状态分析采用在线过滤，并明确区分选择前与选择后时点。
输入若由历史自适应决定，则按已知实验策略条件化，记录其与此前观测的依赖。

\subsection{状态演化算子}
尚未给定实际选择和反馈时，对所有可能观测求和得到
\begin{equation}
\Tt_{\theta,u}=\sum_y\sum_{r\in\Rr}\Ii_{\theta,(y,r)}^{u},\qquad
\langle1,\Tt_{\theta,u}b\rangle=1.
\label{eq:evolution}
\end{equation}
状态演化算子把当前状态分布映射为完成一次决策与学习后的预测分布。
给定预先确定的刺激和反馈协议序列，
\begin{equation}
b_{T+1}^{\mathrm{gen}}=\Tt_{\theta,u_T}\cdots\Tt_{\theta,u_1}b_1.
\end{equation}
自主轨迹模拟通过抽样近似这一传播。
模拟中的反馈必须由模拟回答和实际任务评分规则产生。
自适应输入协议则把输入策略一并放入前向生成过程。

\subsection{各层的用途与分析条件}
学习核 $\Ll$ 描述给定一次内部记录、选择和反馈后，单个认知状态怎样更新。
$\Ii$ 描述观察实际行为后，研究者的状态分布怎样改变；$\Tt$ 描述尚未观察行为时，
所有可能的决策和学习结果形成怎样的状态分布。这些对象由同一单试次核连接。

在同一状态空间和兼容输入协议下，比较 $\Tt_{u_b}\Tt_{u_a}$ 与 $\Tt_{u_a}\Tt_{u_b}$
可定义信息顺序对学习分布的影响。稳态与谱分析还要求明确平稳输入策略、状态空间及有限维近似。
认知状态包含连续的规则信念、判断精度和失败累积量；长期分析需规定这些变量的取值处理及近似误差。
长期谱、稳态及前瞻干预的分析结果尚待计算。

\section{个体参数、数值推断与估计目标}
\label{sec:inference}

\subsection{九项共同参数与条件3的配对参数}
\label{sec:parameters}
条件1、2的个体参数为
\begin{equation}
\theta_s=(p_{0,s},a_{e,s},\lambda_s,g_{0,s},K_s,\delta_s,\beta_{0,s},\eta_s,\kappa_s).
\label{eq:parameters}
\end{equation}
条件3再加入$\gamma_s$。以下省略被试下标；各参数在一次拟合的整段试次内固定，
其共同作用产生随时间变化的状态和行为。

{\small
\begin{longtable}{L{1.3cm}L{3.0cm}L{5.2cm}L{4.7cm}}
\toprule
参数 & 名称与范围 & 在模型中的作用 & 解释时的边界\\
\midrule\endhead
$p_0$ & 基础搜索概率，$(0,1)$ & 失败累积量为零时的搜索概率。 & 实际搜索率还由$a_e$与$F_t$决定。\\
$a_e$ & 失败敏感度，$[0,\infty)$ & 失败量对搜索对数优势的增益。 & 为零时，$\lambda$失去行为作用。\\
$\lambda$ & 失败累积量保留，$[0,1]$ & 完全成功后保留的失败量比例。 & 作用依赖成功与失败的时序。\\
$g_0$ & 全局搜索比例，$[0,1]$ & 搜索发生时选择全局目录的概率。 & 局部路线存在邻域为空时的功能回退。\\
$K$ & 活跃容量，整数$\geq2$ & 同时保留的规则数上限。 & 是当前规则表示下的有效容量。\\
$\delta$ & 规则信念保持，$[0,1]$ & 零分反馈时的信念保持率；本题有效值为$\delta+(1-\delta)r_t$。 & 正确反馈时有效值为1；其作用依赖反馈时序。\\
$\beta_0$ & 初始精度，$[0.1,25]$ & 首题规则的类别距离敏感程度。 & 后续精度还受学习和继承影响。\\
$\eta$ & 精度学习率，$[0,1]$ & 控制$10\eta$倍反馈梯度的更新幅度。 & 为零时，全部规则保持$\beta_0$。\\
$\kappa$ & 执行维持倾向，$[0,\infty)$ & 将当前执行规则的信念乘以$e^\kappa$，再归一化选择下一执行身份。 & 作用取决于规则信念及容量保留；被移除时发生强制切换。\\
$\gamma$ & 配对信念保持，$[0,1]$ & 条件3中各规则的条件配对分布保持率。 & 仅条件3估计；与精度存在补偿可能。\\
\bottomrule
\end{longtable}}

上述含义界定参数在模型中的作用。把点估计进一步解释为稳定个体特质，需要检查：
不同起点能否找到相近的高似然区域、相近分数下参数是否分散、模拟真值能否恢复，
以及关键内部状态是否对参数变化稳健。已知零作用参数在对应参数点单独标记。

当$\eta=0$时，反馈引起的精度增量为零。首题规则的精度为$\beta_0$，
后续新规则继承的平均精度也为$\beta_0$，因此全部活跃规则在整段任务中保持相同精度。
在此边界上，规则信念和活跃集合仍可随反馈变化。

\subsection{粒子滤波：积分可能的内部路径}
第$n$个粒子表示一条潜在路径，作答前权重为$W_t^{(n),\mathrm{pre}}$。
粒子内的执行身份为$e_t^{(n)}$，其选择概率为$p_t^{(n)}(c)=p_{t,e_t^{(n)}}^{(n)}(c)$。
结合该粒子抽取的内部刺激，研究者的作答前预测为
\begin{equation}
P_t(c)=\sum_{n=1}^{N_{\mathrm{PF}}}W_t^{(n),\mathrm{pre}}p_t^{(n)}(c),\qquad
\sum_nW_t^{(n),\mathrm{pre}}=1.
\label{eq:pf-predict}
\end{equation}
实际选择出现后，路径权重更新为
\begin{equation}
W_t^{(n),\mathrm{choice}}=
\frac{W_t^{(n),\mathrm{pre}}p_t^{(n)}(y_t)}{P_t(y_t)}.
\label{eq:pf-update}
\end{equation}
更能预测实际回答的路径获得更高后验权重；所有路径再根据实际反馈完成各自的学习。
有效样本量（effective sample size，ESS）衡量粒子权重的集中程度，其定义为
\begin{equation}
\mathrm{ESS}=\frac{1}{\sum_n\bigl(W_t^{(n),\mathrm{choice}}\bigr)^2}.
\end{equation}
当ESS低于预先设定的阈值时进行重采样，
子粒子继承认知状态，并产生独立的后续随机分支。
粒子数、阈值和重复次数属于数值设置，按具体协议记录。

可在指定时点汇总规则在场概率、规则信念、执行身份和精度分布。
选择前分布使用此前的选择历史；选择后分布再以本题选择条件化，二者均在本题反馈前记录。
当前参数由完整有效记录估计，因此整段状态属于拟合后的回顾性重建。
每题过滤仍只输入到该时点已经发生的选择与反馈；全记录估参的信息用途另行注明。
固定参数下的路径区间反映潜在过程的不确定性；独立计算重复的标准误反映数值误差；
个体参数的不确定性需另行传播。三者在图表中分别注明。

\subsection{完整记录似然：评价实际选择序列}
设被试共有$T$道有效题。参数$\theta$在整次任务中固定，规则集合、信念、精度和执行身份逐题变化。
第$r$次独立粒子滤波在观察本题选择之前计算$P_{t,r}(y_t;\theta)$，
随后用实际选择与反馈推进学习。该次计算的完整序列似然及合并后的每题负对数似然为
\begin{align}
\widehat L_r(\theta)&=\prod_{t=1}^{T}P_{t,r}(y_t;\theta),\\
\widehat{\mathrm{NLL}}(\theta)&=-\frac1T\log\left[\frac1R\sum_{r=1}^{R}\widehat L_r(\theta)\right].
\label{eq:objective}
\end{align}
$R$为独立滤波次数，$T$包括首题和末题。实现使用对数求和，避免长序列概率相乘时下溢。
先合并各次完整序列似然，保持每次计算内部的连续学习历史。
此前逐题平均预测概率再取对数的分数保留为诊断量，两个有限计算结果可能不同。
本题选择影响过滤权重，反馈影响学习者状态；两个更新的含义见第\ref{sec:inference}节前文。

独立滤波是计算误差的重复单位。逐次删除一个完整系统，计算删一法蒙特卡洛标准误
（Monte Carlo standard error，MCSE）。候选之间使用相同种子编号配对，另以新种子确认。
同时记录每个系统的相对序列似然权重$\widehat L_r/\sum_j\widehat L_j$。
少数系统占据大部分权重提示计算不稳定；较小的重复标准误仍需结合粒子数检查理解。
有限重复只能发现部分计算误差，不能保证已消除粒子近似偏差。
每个最终确认点另用新种子进行两次1024粒子过滤。两个粒子数的NLL绝对差，加两倍MCSE平方和的平方根，
须不超过0.005，且高粒子计算的系统最大权重不超过0.9，才通过粒子数检查。
候选及参考点均须通过该检查；敏感、边界和未完成结果分别记录。

\subsection{联合参数候选集合与工作判据}
\label{sec:selection}
当前目标是得到若干完整参数组合，说明哪些组合能同时解释真实选择和自主轨迹形态。
每个候选都计算完整选择似然，并生成多条自主路径。候选比较保留三个层次：
\begin{enumerate}
\item 选择拟合：相对于已独立确认候选中的最低NLL，记录每题差值$\Delta(\theta)$。
\item 数值精度：记录差值和各点NLL的MCSE，以及完整系统的最大权重。
\item 生成相容性：检查被试的单路径形态在该参数生成分布中的位置，并记录全部形态的概率质量。
\end{enumerate}
主工作容忍为每题NLL差值0.01，同时显示0.005与0.02的结果。
若$\Delta+2\mathrm{MCSE}_{\Delta}\leq\epsilon$，则通过容忍为$\epsilon$的选择拟合检查；
若$\Delta-2\mathrm{MCSE}_{\Delta}>\epsilon$，则超出该容忍；其余记为边界未确定。
两个被比较点及差值的两倍MCSE均须不超过0.005，完整系统最大权重须不超过0.9。
两倍MCSE不超过0.002另记为较高数值精度。阈值是本次先导的工作精度，删一法误差带表示有限计算的不确定性。

通过选择检查、数值检查和下述形态相容性检查的完整组合，构成已检验的联合候选集合。
NLL容忍和生成尾部阈值均需后续重复模拟校准，其集合目前没有置信区域或参数后验的含义。
输出保存每一组完整参数；逐个参数的最小值与最大值只概括该集合，不能任意交叉组合。
相近行为评分下可能存在不同内部状态解释，相关解释另作敏感性检查。

若没有组合通过，分别报告：搜索覆盖不足、评分仍有数值误差、选择与生成要求有冲突，或已检查组合的生成形态不相容。
预算内可以检查边界和补充重复；达到预算后交付空集合及原因。有限候选失败所支持的结论，限于本次机制和已检查范围。

\subsection{自主路径：先描述单条形态，再比较分布}
\label{sec:shape}
给定一个参数组合，从模型初始状态开始模拟选择、获得相应反馈，并沿自身经验学习。
同一组合可能产生多种路径。多条路径的均值缓慢上升，可能来自各条路径在不同时间突然改善；
均值水平且范围宽，也可能来自持续较高与持续较低路径的混合。因此，形态指标先在每条完整路径上计算。

本次先导固定被试已记录的刺激日程及记录长度。原实验还按每64题的表现调整后续类别数量，
并根据各类别达标情况提前结束任务。本次条件生成检验覆盖给定日程下的学习形态；
自适应刺激安排和完整实验学习时长分布另需复现实验程序。全部路径保留，未学会及末段随访不足均记录分母。

令$z_t=\ind\{y_t=c_t\}$表示完全正确。条件3同时保留原反馈得分，避免混淆完全正确与部分正确。
为描述形态，在每条正确性序列上拟合四参数曲线
\begin{equation}
q(t)=l+(u-l)\sigmoid\bigl((t-c)/s\bigr),\qquad s>0.
\label{eq:shape-sigmoid}
\end{equation}
$l,u$表示两端水平，$c$表示变化中心，$s$控制变化持续时间。
曲线以二值观察的对数损失拟合，允许$u<l$，因此下降路径也可描述。
其10\%至90\%变化宽度为$w=2\log(9)s$，用于区别短暂变化与长时间变化。
该曲线是行为摘要；原认知模型继续负责选择概率及状态变化。

与整段固定正确率相比，曲线改善的描述性BIC差至少为6、且$|u-l|\geq0.15$时，记录为清楚的水平变化。
BIC差由两倍拟合对数似然增益减去$3\log T$计算，只作为预定的描述门槛。
未达到门槛时，两端水平记为整段正确率，变化宽度记为缺失，并另存变化证据标志。
上升和下降均按实际变化终点计算随访。起点超出记录、变化终点超出记录或终点后不足64题，均记录截断标志。
这些标志防止把观察时间不足解释为已经验证的保持。

每条路径的主要特征包含三组：
\begin{enumerate}
\item 水平与保持：两端水平，以及变化完成后最多64题的正确率；缺少随访时记录实际分母。
\item 变化形式：正向改善幅度、经$\log(1+w)/\log(1+8T)$缩放的持续宽度、变化证据与截断标志。
\item 波动：16、32、64题块正确率的相邻差均方根，以及32题块相对拟合趋势的残差离散程度。
\end{enumerate}
残差离散程度为各块$n_b(\bar z_b-\bar q_b)^2/\max\{\bar q_b(1-\bar q_b),0.02\}$的均值，
其中$n_b$是块内题数。进入距离时，该量截在10并除以10。三组特征在距离平方中分别分配0.35、0.35、0.30权重，
组内按维数平均。绝对变化中心$c$单独保存，不进入主要形态距离。边界截断和有限窗口使该距离只近似容许时间平移。

为了展示路径组成，另按预定规则标记明显波动、无清楚改善、持续较高、下降、截断变化、改善后回落、
短时改善并保持、渐进改善并保持和中间宽度。明显波动优先记录；持续较高门槛为整段0.8，
残差离散程度大于3记为明显波动。保持随访比拟合高水平低0.1以上记为回落。
充分观察到的正向变化，宽度不超过32题记为短时改善，至少64题记为渐进改善，其余记为中间宽度。
所有标签及连续特征一起保存；这些行为形态的心理含义需要额外的内部过程证据。
形态比例报告全部路径计数和模拟二项区间，逐条路径持续较低所占比例也保留。

\subsection{全分布评分与生成相容性}
设$\bm f_m$为第$m$条自主路径的形态特征，$\bm f_{\mathrm{obs}}$为被试特征。
Energy评分同时计入模拟与观察的距离以及模拟分布自身的离散程度：
\begin{equation}
\widehat{\mathcal E}(\theta)=
\frac1M\sum_{m=1}^M\norm{\bm f_m-\bm f_{\mathrm{obs}}}_2
-\frac{1}{2M(M-1)}\sum_{m\ne n}\norm{\bm f_m-\bm f_n}_2.
\label{eq:energy}
\end{equation}
$M$为全部自主路径数。评分越小，表示所选特征下的预测分布评分较好。
固定32题位置的正确率和类别条件反应频率仍计算既有Energy诊断，便于看见时间分布、反应混淆与主要形态检查遗漏的差距。

生成相容性另用独立模拟校准。将确认阶段的256条路径按预定种子顺序分成128条参考路径和128条校准路径。
目标路径到参考集合中最近$\lfloor\sqrt{128}\rfloor=11$条路径的平均欧氏距离为差距量。
多邻居距离检查目标附近的模拟质量，同时避免两群之间的空隙仅因接近均值而通过检查。
比较被试差距与各条校准路径的差距，
记后者大于等于被试差距的次数为$k$，校准条数为$n$；报告尾部秩$(k+1)/(n+1)$及$k/n$的95\% Wilson模拟区间。
区间下界大于0.05记为相容，上界小于0.05记为不相容，跨越0.05记为未确定。
区间衡量固定参考集合下的模拟计数误差。它未包含参数不确定性、参考集合误差或指标选择误差。
主要距离用全部路径计算，单条最近路径及很宽的逐点范围带仅作描述。

形态指标会压缩数据；相容表示通过已规定特征的工作检查。任务时点、类别混淆和原始路径仍随结果展示。
绝对改善时间单独报告分布和截断比例。参数影响这一时间分布，同一参数也允许过程随机性带来较大差异。
频域摘要暂作可选探索，因为上升和时间反转后的下降可以有相同功率谱，其心理解释需保留方向、幅度和保持信息。

\subsection{六名被试先导：范围、预算与停止规则}
2026年10月26日获准实施的首批案例为常规开发案例110、205、306，以及此前生成检验较困难的129、229、317。
三种条件各两人，完整有效题数依次为448、704、768、256、1088、1024。
常规案例依据已有开发样本的记录长度选取；困难案例依据既有生成偏差选取。两组均有选择用途，群体代表性尚待检验。

每人先评价8个起点，包括既有似然与联合选参点、既有确认候选和较分散的粗取值组合。
随后最多4轮，每轮增加4个联合候选，总数最多24。初筛每点两次128粒子滤波，并生成32条路径。
每人至多4个候选使用新种子做四次512粒子滤波、两次1024粒子对照，并各生成256条路径。
仍有数值疑问时，在预算允许下为该人的所有确认点各增加四次512粒子滤波；
完整追加批次才能并入配对比较。部分完成的结果单独保存。
候选同时保留似然较好、形态较好及不同参数组合；确认完成后按各门槛输出集合。

粗取值以较疏的点检查范围：$p_0\in\{0.005,0.05,0.2,0.65,0.95\}$，
$a_e\in\{0,0.5,2,8\}$，$\lambda\in\{0,0.6,0.95,1\}$，$g_0\in\{0,0.2,0.5,1\}$，
$K\in\{2,4,8\}$，$\delta\in\{0,0.5,0.9,1\}$，$\beta_0\in\{0.1,1,5,15,25\}$，
$\eta\in\{0,0.003,0.03,0.3,1\}$，$\kappa\in\{0,1,3,8\}$，条件3的$\gamma\in\{0,0.6,0.95,1\}$。
这些值用于提出联合候选，实际只评价有限组合。后续沿较好候选的相邻取值和中点细化；
若最优点触及$K=8$、$\kappa=8$或$a_e=8$，最后一轮允许一次相应的10、12或12外扩检查。
$a_e=0$时记录$\lambda$无作用，并以统一值去重。较疏取值节约搜索资源，可辨识性仍由行为差异和恢复证据决定。

总预算96 CPU小时：评分校准12、搜索32、独立似然确认36、自主生成12、指标分析与报告4。
使用至多128个可用CPU、每工作进程一个数值线程、单层进程并行；任务不足时按实际任务数运行。
每个数值工作预留CPU并设置进程上限，阶段预算耗尽即停止增加工作。候选上限、重复上限及阶段上限共同构成停止规则。
耗时取决于容量、执行维持、可并行任务和系统负载；既有单次成本提示约5--8小时，正式校准后更新预计时间。

本轮交付六人的完整参数表、三档容忍集合、数值未确定项、全部自主路径及形态概率、失败原因和判断日志。
完整重复参数恢复与状态恢复另行安排；它们分别检验生成真值的参数集合覆盖和潜在状态的可重建程度。
执行配置为\texttt{configs/shared/model1024\_individual\_regions.yaml}，逐次结果保存输入、代码和随机种子的身份信息。
本节记录修订2的冻结方法。首版48次校准和首批96次过滤已完成，共7.118 CPU小时，形态生成尚未开始。
校准发现部分点对粒子数敏感，因此增加最终候选的1024粒子对照，并在96 CPU小时内重新分配预算。
既有144次过滤的机制、评分和种子保持一致，原始结果保留并复用，已用CPU继续计入预算。
计算完成与科学通过按最终结果分别更新。

\section{口述：对当前作答依据的带噪声检验}
\label{sec:oral}

\subsection{检验对象与时点}
以下规定口述检验的设计。当前资料尚待用户最终人工检查，既有结果保留暂定状态；
本轮参数研究使用行为记录与模拟资料，新增口述验证及编码修改暂缓。
口述指导语要求说明：当前刺激的哪些特点使被试选择这个选项。
因此，主要检验对象为\textbf{本题选择之后、反馈之前的局部作答依据}。
模型仅由行为估参和推断状态；实际口述经编码成为一组判断约束，
再检验这些内容与模型推断的作答依据相容到什么程度。

口述是内部表征的近似检验，报告还受到表达、遗漏、参照歧义和事后解释的影响。
口述检验考察报告内容与推断的学习过程是否相容，并保留不同报告层次的噪声与覆盖限制。
口述质量分层使用独立于模型成绩的标准，常规多部位表达的低目录覆盖单独记录。

\Needspace{10\baselineskip}
\subsection{局部规则分量的解释质量}
设规则$h$对所选类别的区域有局部分量$j$，$\alpha_{t,h,j}^{(n)}$为该粒子内的分量权重。
选择到内部刺激最近的分量，并列时均分；分量权重之和为1。
实际选择之后，粒子、规则和分量的联合解释质量，即分配给各潜在作答依据的概率质量，为
\begin{equation}
A_t(n,h,j\mid y_t)=
\frac{W_t^{(n),\mathrm{pre}}\ind[h=e_t^{(n)}]p_{t,h}^{(n)}(y_t)
\alpha_{t,h,j}^{(n)}}{P_t(y_t)}.
\label{eq:oral-responsibility}
\end{equation}
指示函数把每条路径的解释质量赋给其执行规则，其他活跃规则在该路径中的解释质量为0。
该质量对全部$n,h,j$求和为1。它按当前选择对潜在路径条件化一次，
再根据各路径执行规则的局部分量评价口述内容。

对一条报告的完整约束，结构相容要求全部已编码判断在局部规则中获得显式支持；
功能相容要求整个局部分量区域蕴含这些数值约束。
区域与报告约束没有正体积交集时记为冲突，其他信息不足情形记为未确定。
同一报告的多个判断在同一局部分量内联合评价。
把式\eqref{eq:oral-responsibility}中相应质量累加，分别得到
\begin{equation}
M_t^{\mathrm{compatible}}+M_t^{\mathrm{conflict}}+M_t^{\mathrm{undetermined}}=1.
\end{equation}
这些指标量化内容相容性，报告措辞的生成概率需要另外定义表达模型。

\subsection{编码、覆盖与解释边界}
编码保存原话、解析覆盖、参照歧义和目录范围。
带“一半”的长短按长度范围中点解释；无明确参照的“长／短”具有歧义，
保留预定的中点与相对读法作敏感性分析。
在同一报告集合上使用预定读法与判定标准，并报告相容、冲突、未确定质量及覆盖率。

覆盖率检查目录中是否至少存在一个分量能够支持整条报告，可在模型评价之前确定。
目录覆盖低意味着该对齐方法对相应内容的判别力有限。
多部位报告可能受到目录覆盖的限制，因此报告内容与方法覆盖率分别记录。
对于目录可覆盖内容，进一步比较相容规则是否在场、承担多少作答贡献，才能定位状态或机制偏差。

口述未进入个体参数拟合，部分口述已参与机制开发。
开发样本上的吻合支持模型解释的可行性，独立验证需使用事先冻结的方法及独立资料。
状态变化对口述解释的贡献，可通过同一模型的动态表征与静态表征进行比较。

\section{过程预测与模型检验}
\label{sec:validation}

\subsection{规则进入、信念变化与执行切换}
规则信念、规则管理和执行状态共同形成行为轨迹。
信念提高会增加规则成为局部搜索起点的概率、提高容量替换时的保留优先级，
并在其他条件相同时提高它被选为下一执行规则的概率。
执行维持倾向使已有作答方式可以持续，执行切换则改变直接产生选择的规则。

关键过程包括目标规则首次进入活跃集合、获得较高信念、首次执行，以及持续执行的长度。
正确规则进入后，学习者可以继续执行另一条规则；切换到正确规则后，行为可能较快改善。
这构成模型对突然改善的一种过程解释。改善幅度及其保持还取决于知觉、判断精度、反馈学习和后续搜索。

持续执行的规则若系统性预测错误反应，模型可以生成低于均匀猜测基准的正确率。
这种表现的大小受规则目录、刺激类别比例和反应偏好共同影响。
二分类与四分类的均匀猜测正确率分别为0.5和0.25；条件3的均匀猜测平均得分为0.375。
评价低表现时同时记录类别及反应比例，并计算二者独立时的期望正确率。
突变和低表现的筛选标准、窗口长度及评价样本在分析前固定，开发筛选与独立验证分别记录。

\subsection{三类错误来源的定义与推断}
\label{sec:error-sources}
令$\mathcal C^\star\subseteq\Hh$为正确规则集合。当前检验将训练刺激上给出的类别标签
全部符合任务分类的规则纳入该集合；这一操作定义以既定刺激和规则目录为条件。
令$c_t^\star$为本题真实类别，错误选择满足$y_t\ne c_t^\star$。
条件3的半分与零分均属于种类判断错误。

在一条具有明确执行身份的路径中，每次错误属于以下三个互斥且完备的事件之一：
\begin{align}
E_{1,t}&=\{\A_t\cap\mathcal C^\star=\varnothing\},\label{eq:error-absent}\\
E_{2,t}&=\{\A_t\cap\mathcal C^\star\ne\varnothing,\ e_t\notin\mathcal C^\star\},\label{eq:error-other}\\
E_{3,t}&=\{e_t\in\mathcal C^\star\}.\label{eq:error-executed}
\end{align}
第一类是正确规则当前缺席；第二类是正确规则在场，但本题执行其他规则；
第三类是已经执行正确规则，仍因知觉和类别概率选择产生错误。
这里的“在场”对应当前活跃集合，曾经进入后又被移除的规则在当前题属于缺席。

实际行为中，研究者对这些事件的判断具有不确定性。
对一次已观察到的错误选择，使用式\eqref{eq:pf-update}的选择后、反馈前权重计算
\begin{equation}
\widehat{\Pr}(E_{j,t}\mid y_{1:t},r_{1:t-1},u_{1:t},\widehat\theta)
=\sum_n W_t^{(n),\mathrm{choice}}\ind[E_{j,t}^{(n)}],
\qquad j=1,2,3.
\label{eq:error-posterior}
\end{equation}
三个概率之和为1。它们是在固定参数和模型假设下，对一次错误的来源分配的概率。
规则信念表示路径内学习者的判断；这些事件概率表示研究者根据行为对路径的推断。

\subsection{较高信念与执行不一致的联合事件}
\label{sec:lag}
在场但未执行还可进一步区分：正确规则的信念是否已经达到较高水平。
将正确规则的总信念记为$b_t^\star=\sum_{h\in\A_t\cap\mathcal C^\star}\pi_t(h)$，
当前分析用下列联合事件描述较高信念与其他规则执行同时出现：
\begin{equation}
L_t=\{\A_t\cap\mathcal C^\star\ne\varnothing,\ b_t^\star\geq0.5,
\ e_t\notin\mathcal C^\star\}.
\label{eq:lag}
\end{equation}
其后验概率按每条路径上的联合事件加权计算。
该事件在答对或答错时均可发生；另一条规则也可能在某个刺激上产生正确选择。
阈值0.5包含信念并列情形，且零维持参数下的执行抽样也可能产生该事件。
因此，事件的出现频率、连续长度及其对$\kappa$的依赖分别评价。

状态图分别记录正确规则在场、较高信念和实际执行的概率。
某个边际概率首次达到0.5的试次定义为检测时点，它随推断方法和阈值变化。
真实生成路径的首次事件时点则由路径直接给出；首次事件时点的后验需要保留完整路径顺序。
事件未发生、发生但未检出，以及可评价试次数分别记录。

\subsection{行为比较、参数恢复与状态恢复}
行为比较使用共同的数据用途、知觉设置和规则目录，为各候选机制分别估计参数。
本轮完整记录先导输出联合候选集合，选择与生成检查均参与选参，结果按开发解释评价。
比较冻结后的选择似然、自主学习轨迹及模型复杂度，检查较好的预测是否在独立数值重复中保持。
执行维持的作用结合零维持边界检验，并记录参数边界、相近似然点及不同起点的结果。

参数恢复在已知参数下生成完整选择与反馈序列，再仅根据生成行为重新估计参数。
它检查执行维持与信念保持、精度、搜索及容量之间的补偿关系。
恢复检查生成点是否落入重新估计的联合候选集合、集合大小及其行为预测差异。
把重新估计点的完整记录似然与生成点比较，可检查搜索是否遗漏较好候选；
同一参数多次生成的集合覆盖需要重复模拟评价。本轮六人先导将完整重复恢复另列预算。
当前有限数据和有限搜索下的恢复结果，约束相应条件下的心理解释。

状态恢复保留生成路径的活跃规则、信念、精度、配对信念和执行身份，
再从行为重建这些状态。分别使用生成参数和重新估计的参数重建，
以评价参数估计对状态推断的影响；两种重建均记录有限粒子误差。
三类错误逐类统计真实次数、正确分类数、误报与漏报，避免多数类主导总体正确率。

联合事件的恢复同时评价概率与分类。Brier误差是预测事件概率与真实事件指示值之间的
平均平方误差，越小表示概率预测越准确；对数评分评价分配给真实事件的概率。
当前事件分类阈值为0.5，高置信误报阈值为0.8。
首次检测时点、持续长度及真实事件覆盖范围与上述评分一起报告。

\subsection{当前检验范围与解释边界}
规则目录规定模型能够表达的判断形式，固定类别标签规定各类别区域对应的反应。
因此，规则发现、容量、正确规则集合和口述相容性的解释均以该目录为基础。
口述内容超出目录表达范围时，覆盖率用于说明相关评价的判别能力。

反馈调节保持率与入场份额规定成功、失败及部分成功如何改变已有判断和新规则的试用程度。
精度与配对继承规定任务知识如何延续到新规则，执行维持倾向规定已有作答方式如何延续。
这些心理假设的适用性由条件预测、自主生成与恢复检验共同评价。

截至2026年10月24日的执行机制开发检验包含九名被试，以及30个用于机制和参数恢复的模拟案例。
自主生成的后段平均正确率在九名被试上均低于真实值，明显改善及其保持仍有待充分刻画。
在其中24条具有离散执行身份的模拟路径中，第二类错误共出现60次；
重新估计参数后的状态分类正确识别0次，使用各路径的生成机制及真实参数重建时正确识别2次。
这些分类每次取三类错误中后验概率最高的一类。较高信念而执行其他规则的事件也较少，
当前模拟对长期执行滞后的覆盖有限。

上述结果表明，当前推断对“正确规则在场但未执行”的逐题识别仍较弱。
恢复中也发现部分重新估计点的似然低于生成点，说明优化过程限制了当前恢复结果。
本模型保留为研究候选；真人状态图按模型条件下的概率描述解释。
个体参数和内部状态的稳定心理解释仍需更充分的恢复与独立资料检验。
状态估计目前以冻结参数为条件，参数估计的不确定性尚未积分到状态分布中。

2026年10月25日至26日的追加研究已完成共同选参、定向续跑和状态推断检查。
共同选参的27组中9组通过独立确认：9组后段NLL改善，6组自主分布改善，317的3组自主分布变差。
两参数条件恢复在12条固定路径中精确恢复2条，多个组合具有相近分数。
最新状态阶段新增48条完整模拟轨迹并完成108次过滤；既有明确逐题标签尚未通过完整验收，
部分低概率范围和整段在场量获得有限支持。少数完整历史系统检查的数值表现存在案例差异。
这些结果支持优先核查个体参数集合、数值稳定性和生成形态，持续执行仍保留为研究候选。
本轮六人先导沿用这一证据边界，逐项结果以研究指南和对应报告为依据。

\appendix
\section{固定设置与初始化}
\label{app:settings}

下表汇总模型的固定设置。个体参数的定义与范围见第\ref{sec:parameters}节。

{\small\noindent
\begin{tabularx}{\linewidth}{L{3.3cm}Y}
\toprule
项目 & 规定\\
\midrule
基础规则权重 & 目录中每条规则的基础权重为$\mu(h)=1/J$。\\
首题状态 & 从基础权重抽取一条规则并执行它；信念为1，判断精度为$\beta_0$，失败累积量为0。\\
首题配对知识 & 条件3的三种配对等概率，各为$1/3$。\\
知觉误差 & 分布形式与参数由独立知觉测量确定，在行为参数估计期间固定。\\
判断精度 & 取值限制为$[0.1,25]$，反馈梯度的更新步长为$10\eta$。\\
额外随机作答率 & 固定为0；每条路径的选择概率由当前执行规则给出。\\
局部结构搜索 & 在非空结构邻域中，按编辑距离的负指数归一化抽取候选。\\
局部功能搜索 & 结构邻域为空时使用功能距离，指数尺度为$0.10$；距离以100000个均匀参考刺激估计，随机种子为0。\\
规则信念保持 & 有效保持率为$\delta_t=\delta+(1-\delta)r_t$；幂变换并归一化后结合反馈似然。\\
容量替换 & 新规则替换更新后信念最低的已有规则；并列时等概率选择，执行规则适用同一规定。\\
新规则信念 & 有空位时获得$1/(|\A_t|+1)$；满容量时获得$r_t\pi_t^+(h^-)+(1-r_t)/K$。幸存规则按原比例分配剩余质量。\\
新规则继承 & 判断精度继承替换前全部活跃规则的信念加权平均值；条件3继承综合配对知识。\\
执行准备 & 学习与规则管理后，按规则信念及执行维持倾向抽取下一身份；下一题呈现刺激前确定。\\
\bottomrule
\end{tabularx}}

\section{规则目录}
\label{app:catalogue}

\subsection{二分类目录：29条}
{\small\noindent
\begin{tabularx}{\linewidth}{L{4.4cm}Yr}
\toprule
规则形式 & 判断例子或边界 & 数量\\
\midrule
单特征阈值 & $x_i\leq0.5$ & 4\\
两个特征的顺序 & $x_i\leq x_j$ & 6\\
两个特征的和 & $x_i+x_j\leq1$ & 6\\
两组特征和的顺序 & $x_i+x_j\leq x_k+x_l$ & 3\\
两个特征是否接近 & $|x_i-x_j|\leq0.10$ & 6\\
单特征是否在中间 & $|x_i-0.5|\leq0.10$ & 4\\
\midrule
合计 & 类别区域与标签固定 & 29\\
\bottomrule
\end{tabularx}}

\Needspace{18\baselineskip}
\subsection{四分类目录：152条}
{\small
\begin{longtable}{L{2.6cm}L{11.1cm}r}
\toprule
涉及特征数 & 规则形式 & 数量\\
\midrule\endhead
两个特征 & 两个单特征阈值交叉 & 6\\
两个特征 & 同一特征对的顺序与和交叉 & 6\\
三个特征 & 阈值与另两特征顺序 & 12\\
三个特征 & 阈值与另两特征和 & 12\\
三个特征 & 三阈值条件分支 & 24\\
三个特征 & 阈值分支下顺序或和 & 24\\
四个特征 & 不重叠特征对的顺序交叉 & 3\\
四个特征 & 不重叠特征对的和交叉 & 3\\
四个特征 & 顺序分支与剩余阈值 & 12\\
四个特征 & 和分支与剩余阈值 & 12\\
四个特征 & 最大或最小特征分类 & 2\\
两个特征 & 阈值与共享特征顺序 & 12\\
三个特征 & 围绕同一部位两次比较 & 12\\
三个特征 & 两部位接近与第三部位长短交叉 & 12\\
\midrule
合计 & 固定标签下的规则判断与类别区域 & 152\\
\bottomrule
\end{longtable}}

\clearpage
\section{主要符号}
\label{app:symbols}

{\small
\begin{longtable}{L{3.0cm}L{11.5cm}}
\toprule
符号 & 含义\\
\midrule\endhead
$S_t,\ m_t$ & 决策就绪状态与本题内部记录。\\
$e_t,\ \kappa_s,\ \omega_{t+1}(h)$ & 本题执行身份、被试层面的执行维持倾向与下一执行身份的选择概率。\\
$\A_t,\ \Hh,\ K$ & 活跃规则集合、固定规则目录与容量。\\

$\pi_t(h),\ \bar\pi_t(h),\ \pi_t^+(h)$ & 作答前规则信念、保持处理后的中间信念、反馈更新后的规则信念。\\
$\delta,\ \delta_t$ & 基础规则信念保持率与由本题反馈决定的有效保持率。\\
$\nu_{t+1},\ \A_t^{\mathrm{keep}}$ & 新规则的入场信念份额与本次规则管理后的幸存集合。\\
$\beta_t(h),\ \beta_t^+(h),\ \bar\beta_t$ & 作答时精度、反馈学习后精度与新规则继承的平均精度。\\
$F_t$ & 决策就绪状态中的失败累积量。\\
$q_t,\ \bar q_t,\ q_t^+,\ \gamma$ & 条件配对信念、保持后分布、反馈后分布与配对保持率。\\
$\Pi_t,\ \Pi_t^+,\ Q_t^+$ & 作答前及反馈后的规则与配对联合信念、反馈后的整体配对边际。\\
$\ell_t(h),\ \ell_t(h,k)$ & 规则的反馈似然与给定规则和配对时的反馈似然。\\
$p_{t,h}(c),\ p_t(c)$ & 单条规则的类别预测与单条路径按执行规则产生的选择预测。\\
$P_t(c),\ W_t^{(n)}$ & 研究者跨路径的预测概率与粒子权重。\\
$\Dd,\ \Ll,\ k_\theta$ & 决策核、学习核与单试次联合生成核。\\
$\Ii,\ \Tt$ & 轨迹推断算子与状态演化算子。\\
$\phi,\ \widehat{\mathcal E}$ & 自主路径的行为特征向量及其Energy评分。\\
$\Ee_u$ & 给定任务输入时，由刺激和选择决定的环境反馈分布。\\
$\mathcal C^\star,\ b_t^\star$ & 正确规则集合及其在一条路径内的总信念。\\
$E_{1,t},\ E_{2,t},\ E_{3,t},\ L_t$ & 三类错误的状态事件，以及较高信念与执行不一致的联合事件。\\

\bottomrule
\end{longtable}}
\end{document}
